\section{Deterministic revealing optimization: the exact curve}
\label{sec:optional}

The goal of this section is to prove \cref{thm:optional-curve}, the rigorous
deterministic part of the informal revealing-optimization curve theorem
\cref{thm:intro-common-upper-curves}.  We now allow raw execution for a
finite common time $u$.  All jobs have this same known upper limit, and a
test reveals $p_j\in[0,u]$.  By
\cref{eq:optional-effective-length,lem:offline}, the offline effective length
is $\lambda_u(p_j)=\min\{u,1+p_j\}$, and the offline optimum schedules these
effective lengths in nondecreasing order.  We write $\RdetRO(u)$ for the optimal
deterministic size-asymptotic ratio at this fixed common upper limit.

The answer has six branches.  We first define their transition points.  Put
\begin{equation}
 \varphi=\frac{1+\sqrt5}{2},
 \qquad
 \udiam=\frac{1+\sqrt{3+2\sqrt5}}2
 =1.866760399173862\ldots .
 \label{eq:opt:phi-udiamond}
\end{equation}
Let $s_0>1$ be the root of
\begin{equation}
 s_0^3=(s_0+1)^2,
 \qquad
 \uzero=1+s_0=3.147899035704787\ldots .
 \label{eq:opt:u0}
\end{equation}
For $s=u-1$, put
\begin{equation}
 T_s(z)=sz^2+(s^3+s^2+1)z-(s^2+s+1)(s+2),
 \label{eq:opt:Ts}
\end{equation}
and let $\Rquad(u)$ be the positive root of
\begin{equation}
sR^2+(s^3+s^2+1)R-(s^2+s+1)(s+2)=0.
\label{eq:opt:Rquad-equation}
\end{equation}
Denote the left-hand side by $T_s(R)$.
Equivalently,
\begin{equation}
 \Rquad(u)=
 \frac{-1-u+2u^2-u^3+
 \sqrt{u^6-4u^5+10u^4-6u^3-3u^2+6u-3}}
 {2(u-1)}.
 \label{eq:opt:Rquad-explicit}
\end{equation}

At the other end, let $\zstar>1$ solve $\zstar-3=\ln \zstar$ and define
\begin{equation}
 r=\frac2{\zstar-1},
 \qquad
 \RdetOT=1+r=1.570574096649395\ldots,
 \qquad
 \zstar=4.505241495792883\ldots .
 \label{eq:opt:obligatory-constant}
\end{equation}
For $\varphi+2\le u\le \zstar$, define $\Rmix(u)=1+c$ by the unique
solution $(c,m)$ of
\begin{equation}
 \operatorname{artanh}m-m
 =1-\frac1c+\frac12\ln\!\left(1+\frac2c\right),
 \qquad
 u=1+\frac2c-m,
 \label{eq:opt:mixed-parametrization}
\end{equation}
with $r\le c\le1/\varphi$ and $0\le m\le1/\varphi$.

\begin{theorem}[Exact revealing-optimization curve]
\label{thm:optional-curve}
For every fixed $0<u<\infty$, the optimal deterministic size-asymptotic ratio in
the revealing-optimization model is
\begin{equation}
\boxed{
\RdetRO(u)=
\begin{cases}
1, &0<u\le1,\\[1mm]
u, &1\le u\le \udiam,\\[1mm]
\Rquad(u), &\udiam\le u\le \uzero,\\[1mm]
\displaystyle 1+\frac1{\sqrt{u-1}},
   &\uzero\le u\le\varphi+2,\\[3mm]
\Rmix(u), &\varphi+2\le u\le \zstar,\\[1mm]
\RdetOT, &u\ge \zstar.
\end{cases}}
\label{eq:opt:full-curve}
\end{equation}
For every branch there is an explicit deterministic algorithm satisfying
\[
 \cost(\mathcal A,I)\le \RdetRO(u)\cost(\OPT,I)+O_u(n).
\]
On $u\ge \zstar$ the remainder is one absolute $O(n)$ term, uniform over all
finite $u$.
Conversely, for every deterministic algorithm $\mathcal A$ and every
$\eps>0$, there are arbitrarily large fixed instances $I$ such that
\[
 \cost(\mathcal A,I)
 \ge\bigl(\RdetRO(u)-\eps\bigr)\cost(\OPT,I).
\]
For $u\ge \zstar$, these lower-bound instances may be chosen cap-free:
$p_j<u-1$ for every job.  Thus both witnesses on the plateau are uniform in
the sense needed to pass separately to the obligatory endpoint.
\end{theorem}

The proof combines two kinds of lower bounds.  Binary hidden-stopping
instances give the first four branches, while a continuous capped version
of the harmonic construction gives the mixed branch and the plateau.  Three
explicit policies match them: raw execution, a forced-prefix uniform
threshold, and a history-dependent adaptive threshold.
The upper-bound witness can be read from the following policy map.  The
algorithms themselves are specified formally in
\cref{alg:raw,alg:ute,alg:adaptive-threshold}.
\begin{center}
\begin{tabular}{@{}lll@{}}
\toprule
range & policy & parameter choice \\
\midrule
$0<u\le \udiam$
 & \textsc{Raw} & none \\
$\udiam<u\le \uzero$
 & \textsc{ForcedPrefixUTE} & $b=b(u)$ from \eqref{eq:opt:ute-parameters} \\
$\uzero<u\le\varphi+2$
 & \textsc{ForcedPrefixUTE} & $b=0$ \\
$\varphi+2<u<\zstar$
 & \textsc{AdaptiveThreshold} & $c=c(u)$ from \eqref{eq:opt:mixed-parametrization} \\
$\zstar\le u<\infty$
 & \textsc{AdaptiveThreshold} & $c=r$ \\
\bottomrule
\end{tabular}
\end{center}

\subsection{Policy guarantees}
\label{sec:det-policy-guarantees}

\subsubsection{Low and intermediate caps}
\label{sec:opt:low-upper}

The first algorithm is formalized in \cref{alg:raw}.

\begin{algorithm}[H]
\caption{\textsc{Raw}$(J,u)$}
\label{alg:raw}
\begin{algorithmic}
\Statex \textbf{Input:} jobs $J$ and a finite common upper limit $u$.
\State Fix an arbitrary order $j_1,\ldots,j_n$.
\State \textbf{For} $i=1,\ldots,n$, execute $j_i$ untested for $u$ time.
\end{algorithmic}
\end{algorithm}

For $u\le \uDet{2}$, use \textsc{Raw}.  It is optimal when $u\le\uDet{1}$;
when $u>1$, every effective length $\min\{u,1+p_j\}$ from
\cref{eq:model-effective-length} is at least one, so the ratio is at most
$u$.

For the next interval, let $R=\Rquad(u)$, $s=u-1$, and define
\begin{equation}
 K=s^2+s+1,
 \quad D=K+sR,
 \quad b=\frac{K-s^2R}{D},
 \quad \chi=\frac{1-b}{R},
 \quad \alpha=\frac KD.
 \label{eq:opt:ute-parameters}
\end{equation}
The defining equation for $R$ gives
\begin{equation}
 0\le b<\chi<\alpha<1,
 \qquad
 \chi=\frac{s(s+1)}D,
 \qquad
 b=(s+1)\alpha-s;
 \label{eq:opt:ute-order}
\end{equation}
For completeness, these inequalities follow without numerical root
isolation.  Since $T_s$ is increasing on positive arguments,
\begin{equation}
 T_s\!\left(\frac K{s^2}\right)
 =\frac K{s^3}\bigl(-s^3+s^2+2s+1\bigr)\ge0
 \qquad(s\le s_0),
\end{equation}
and hence $b=(K-s^2R)/D\ge0$, with equality only at $s=s_0$.  Moreover,
\begin{equation}
 \alpha-\chi=\frac1D,
 \qquad
 \chi-b=\frac{s^2R-1}{D}>0.
\end{equation}
For the last sign,
\begin{equation}
 9T_s(5/3)=f(s):=6s^3-12s^2-2s-3<0.
 \label{eq:opt:T-five-thirds}
\end{equation}
throughout the present interval.  Indeed, on $[4/5,1]$ the polynomial is
decreasing and $f(4/5)=-1151/125$; on $[1,s_0]$ it has one interior minimum,
while $f(1)=-11$ and $f(s_0)<f(13/6)=-95/36$.  Here
$s_0<13/6$ follows by substituting $13/6$ in
$s^3-(s+1)^2$.  Thus $R>5/3$, while $s>4/5$, and
$s^2R>16/15$.  Finally $K<D$ gives $\alpha<1$.  In particular, $b=0$
occurs only at $u=\uDet{3}$.
The remaining two elementary branches use the single procedure in
\cref{alg:ute}.

\begin{algorithm}[H]
\caption{\textsc{ForcedPrefixUTE}$(J,u,b)$}
\label{alg:ute}
\begin{algorithmic}
\Statex \textbf{Input:} jobs $J$, finite $u>1$, and $b\in[0,1]$.
\State Fix a test order $j_1,\ldots,j_n$; set
       $k\gets\lfloor bn\rfloor$, $\theta_u\gets\min\{1,u-1\}$,
       and $Q\gets\varnothing$.
\State \textbf{For} $i=1,\ldots,n$ \textbf{do}:
\State \quad Test $j_i$ for one unit and observe $p_{j_i}$.
\State \quad \textbf{If} $i\le k$ or $p_{j_i}\le\theta_u$, process $j_i$
       immediately.
\State \quad \textbf{Otherwise}, add $j_i$ to $Q$.
\State Process the jobs of $Q$ in nondecreasing revealed processing time.
\end{algorithmic}
\end{algorithm}

On $\uDet{2}\le u\le \uDet{3}$, use
\textsc{ForcedPrefixUTE} with $b$ from \cref{eq:opt:ute-parameters}.  Rounding
$bn$ changes the cost by $O_u(n)$.

Here is the complete binary calculation.  Adjacent exchanges place as many
long jobs as possible in the forced prefix and put all suffix-long jobs
before suffix zeros.  If $x$ is the total long mass, then the prefix-long
mass is $a=\min\{b,x\}$ and the deferred-long mass is
$d=(x-b)_+$.  After division by $n^2$, the leading excess
$G(x)=A-R O$ is
\begin{equation}
G(x)=
\begin{cases}
\displaystyle
 \frac{1-R}{2}+u\left(x-\frac{x^2}{2}\right)
 -\frac{R s}{2}x^2,
 &0\le x\le b,\\[3mm]
\displaystyle
 \frac{1-R}{2}
 +u\left(b-\frac{b^2}{2}\right)
 +(1-b)(x-b)+\frac{s}{2}(x-b)^2
 -\frac{R s}{2}x^2,
 &b\le x\le1.
\end{cases}
\label{eq:opt:binary-upper-gap}
\end{equation}
The second quadratic is strictly concave, has its stationary point at
$x=\alpha$, and satisfies $G(\alpha)=0$; after substituting
\cref{eq:opt:ute-parameters}, the latter identity is exactly
\cref{eq:opt:Rquad-equation}.  Since $b<\alpha$, its right derivative at $b$
is positive.  The left derivative satisfies
$G'_-(b)=G'_+(b)+s(1-b)>0$, and the first quadratic is concave, so it is
increasing on $[0,b]$.  Thus $G(x)\le0$ for every $x\in[0,1]$, including
the limiting case $b=0$.  This proves the binary upper bound without a
hidden local-versus-global step.

For $\uDet{2}\le u\le2$, arbitrary values reduce to binary values by
independent Bernoulli extremalization, in the following order.  First raise
\emph{every} suffix-deferred value to $u$.  This cannot lower ALG, while its
offline effective length was already $u$.  Only after this deterministic
step, replace a prefix value $p$ by $u$ with probability $p/u$ and by zero
otherwise, and replace a suffix-immediate value $p\le u-1$ by $u$ with
probability $p$ and by zero otherwise.  Notice that the $u$ realization
of a suffix-immediate job becomes deferred, so status is not preserved.
Nevertheless the online check is pairwise.  For jobs $i$ before $j$, the
leading pair contribution is
\begin{equation}
 g_{ij}=\begin{cases}
 1+p_i,&i\text{ is immediate},\\
 2+p_j,&i\text{ is deferred and }j\text{ is immediate},\\
 2+\min\{p_i,p_j\},&i,j\text{ are deferred}.
 \end{cases}
 \label{eq:opt:pair-contribution}
\end{equation}
A rounded prefix job stays immediate and preserves its expected value.  If
$i$ is suffix-immediate, its rounded contribution is one when it becomes
zero and at least two when it becomes $u$, hence its expectation is at least
$1+p_i$.  If $i$ is deferred and $j$ is suffix-immediate, then $i$ has been
raised to $u$ and the expected pair contribution is
$2+up_j\ge2+p_j$; a deferred--deferred pair can only increase.  Thus
expected ALG does not decrease, up to its $O_u(n)$ diagonal terms.

The expected effective lengths are,
respectively,
\begin{equation}
 1+\frac{u-1}{u}p\le \min\{u,1+p\},
 \qquad
 1+(u-1)p\le1+p,
\end{equation}
and $u$ for a deferred value.  Since the minimum is jointly concave,
$\EE\min(\lambda_i',\lambda_j')\le\min(\lambda_i,\lambda_j)$.  Thus some binary realization
has competitive excess at least that of the original instance.

It remains to justify arbitrary values for $2\le u\le \uDet{3}$.  The following
quadratic lemma records the entire endpoint calculation.

\begin{lemma}[UTE endpoint game]
\label{lem:opt:ute-endpoint-game}
For $2\le u\le \uDet{3}$, UTE with the parameters
\cref{eq:opt:ute-parameters} satisfies
\begin{equation}
 \ALG\le\Rquad(u)\OPT+O_u(n).
\end{equation}
\end{lemma}

\begin{proof}
Fix the symbolic prefix/immediate/deferred decisions.  Convexity in one
coordinate reduces a prefix job to $p\in\{0,u\}$, a suffix-immediate job to
$p\in\{0,1\}$, and a suffix-deferred job to $p\in\{1^+,u\}$.  For the last
claim, on $1<p<u-1$ the deferred--deferred term has coefficient
$1-R<0$ multiplying a concave minimum; on $[u-1,u]$, OPT is flat and ALG
is nondecreasing.  Adjacent exchanges give the canonical endpoint order.

Normalize counts by $n$.  Let $a$ be prefix-$u$ mass, $d$ suffix-$u$ mass,
$t$ suffix-$1^+$ mass, and $m$ suffix-immediate-one mass.  Their constraints
are
\begin{equation}
 0\le a\le b,
 \qquad d,t,m\ge0,
 \qquad d+t+m\le1-b.
\end{equation}
Pairwise accounting using \cref{eq:opt:pair-contribution} gives
\begin{align}
 A={}&\frac12+(s+1)\left(a-\frac{a^2}{2}\right)
 +(1-b)(d+t+m)-\frac{m^2}{2}+\frac{s}{2}d^2,
 \label{eq:opt:ute-A}\\
 O={}&\frac12+\frac{(a+d+t+m)^2}{2}
 +\frac{s-1}{2}(a+d)^2.
 \label{eq:opt:ute-O}
\end{align}
Moving $m$ to $t$ does not change $O$ and increases $A$, so $m=0$.  Put
$x=a+d$ and write $F(x)=A-R O$ after optimizing the other variables.
The derivative of the gap in $t$ is
$1-b-R(x+t)$, so its maximizing value is $(\chi-x)_+$.  At fixed
$x,t$, moving mass from $d$ to $a$ leaves $O$ unchanged and changes $A$ at
rate $s(1-x)+b-a\ge0$.  Hence
\begin{equation}
 t=(\chi-x)_+,
 \qquad
 a=\min\{b,x\},
 \qquad d=(x-b)_+.
 \label{eq:opt:ute-reduction}
\end{equation}

For $x\ge\chi$, the gap $A-R O$ is the binary concave quadratic already
discussed; it is tangent to zero at $x=\alpha$.  For $x<\chi$, put
\begin{equation}
 C_0=-\frac{R-1}{2}+\frac{(1-b)^2}{2R},
 \qquad
 \gamma=s-R(s-1).
\end{equation}
Substitution gives
\begin{equation}
 A-R O
 =C_0+a(s+b-sx)-\frac{a^2}{2}+\frac{\gamma x^2}{2}.
 \label{eq:opt:ute-positive-t}
\end{equation}
One has $\gamma>0$ on $1\le s\le s_0$.  Indeed,
$T_s(2)=s(s^2-s+1)>0$, so for $s\le2$ the equation $T_s(R)=0$ gives
$R<2$, and hence $\gamma>2-s\ge0$.  For
$2\le s\le s_0<13/6$,
\begin{equation}
 16T_s(7/4)=12s^3-20s^2+s-4>0;
\end{equation}
the polynomial is $14$ at $s=2$ and has positive derivative thereafter.
Thus $R<7/4$ and $\gamma>7/4-3s/4>1/8$.

On $0\le x\le b$, one has $a=x$, and
\begin{equation}
 A-R O=C_0+(s+b)x
 -\frac{s+1+R(s-1)}2x^2.
 \label{eq:opt:ute-small-x}
\end{equation}
Its derivative is decreasing, and the derivative at $b$ multiplied by $D$
is
\begin{equation}
 R\bigl(1+s^2+R s^2(s-1)\bigr)>0,
\end{equation}
so the maximum is at $b$.  On $b\le x\le\chi$,
\cref{eq:opt:ute-positive-t} is convex and its maximum is at $b$ or $\chi$.
The endpoint $\chi$ lies on the binary branch.  Finally,
\begin{equation}
 F(b)-F(\chi)
 =\frac{\chi-b}{2}\,[2sb-\gamma(b+\chi)],
\end{equation}
and direct use of \cref{eq:opt:ute-parameters} gives
\begin{equation}
 D[\gamma(b+\chi)-2sb]
 =-s+(1+s-s^3)R+s^2(s-1)R^2>0.
 \label{eq:opt:ute-small-certificate}
\end{equation}
For a rational sign certificate, $T_s$ is increasing on positive arguments
and
\begin{equation}
 9T_s(5/3)=6s^3-12s^2-2s-3<0,
\end{equation}
so $R>5/3$.  If the right-hand side of
\cref{eq:opt:ute-small-certificate} is denoted by $q_s(R)$, then
\begin{equation}
 3\,\partial_R q_s(5/3)=7s^3-10s^2+3s+3>0.
\end{equation}
The last polynomial is increasing for $s\ge1$, so $q_s$ is increasing for
$R\ge5/3$.  Finally,
\begin{equation}
 9q_s(5/3)=10s^3-25s^2+6s+15>0.
\end{equation}
For the last sign, write $s=3/2+z$; the expression becomes
$3/2-3z/2+20z^2+10z^3$.  It is at least $3/2$ for
$-1/2\le z\le0$, and for $z\ge0$ it is larger than
$3/2-9/320$.  These inequalities prove that every endpoint gap is
nonpositive.  Diagonals and prefix rounding contribute $O_u(n)$.
\end{proof}
At $u=\uDet{3}$, $b=0$.  The same zero-prefix algorithm remains exact through
$u=\uDet{4}$ and continues with a $\varphi$ guarantee afterwards.  We state
the continuation because it makes the three-algorithm lower envelope in
\cref{fig:intro-curves-det} literal rather than schematic.

\begin{lemma}[Zero-prefix extension]
\label{lem:opt:zero-prefix}
For every $u\ge \uDet{3}$,
\textnormal{\textsc{ForcedPrefixUTE}}$(J,u,0)$ from \cref{alg:ute} is
\begin{equation}
 \max\left\{\varphi,1+\frac1{\sqrt{u-1}}\right\}
 \label{eq:opt:zero-prefix-full-guarantee}
\end{equation}
competitive up to $O_u(n)$.  Thus its guarantee is
$1+1/\sqrt{u-1}$ on $[\uDet{3},\uDet{4}]$ and $\varphi$ on
$[\uDet{4},\infty)$.
\end{lemma}

\begin{proof}
The same coordinatewise reduction leaves capped-deferred mass $d$,
boundary-deferred mass $t$, immediate-one mass $m$, and zeros.  With
$s=u-1$, these masses satisfy $d,t,m\ge0$ and $d+t+m\le1$, and the leading
costs are
\begin{equation}
 A=\frac12+d+t+m-\frac{m^2}{2}+\frac{s d^2}{2},
 \qquad
 O=\frac12+\frac{(d+t+m)^2}{2}+\frac{(s-1)d^2}{2}.
\end{equation}
Move $m$ to $t$ and put $y=d+t$.  For fixed $y$, the ratio is
fractional-linear in $d^2$, so its maximum is on one of two faces:
\begin{equation}
 \frac{1+2y}{1+y^2}\le\varphi
 \quad(d=0),
 \qquad
 1+\frac{2y}{1+s y^2}\le1+\frac1{\sqrt s}
 \quad(d=y).
\end{equation}
The first one-variable function is maximized at $y=1/\varphi$ and the second
at $y=1/\sqrt s$, which gives the two displayed constants.
The second value dominates the first exactly when
$u\le\uDet{4}$.  Binary instances attain the second value.
\end{proof}


\subsubsection{High caps}
\label{sec:optional-upper}

We now prove the upper bounds on the two ranges on which capped outcomes
matter.  Besides completing the upper-bound side of the revealing-optimization curve,
the proof records explicitly why an outcome at the raw cap cannot be
treated as an ordinary deferred outcome.  Throughout this subsection,
$R=1+c$, and we retain $\varphi$, $\uDet{5}$, and $r$ from
\cref{eq:opt:u1-u2,eq:opt:obligatory-constant}.

For every $c>0$ define the history-dependent threshold
\begin{equation}
 a_c(y)=
 \begin{cases}
  1, & y\le-1,\\[1mm]
  \displaystyle
  1+\frac12\ln\!\left(\frac{c+2}{c-2y}\right),
      & -1\le y\le0.
 \end{cases}
 \label{eq:opt:threshold}
\end{equation}
The single algorithm used throughout the last two branches is the following.
It does not use raw execution and therefore does not depend on $u$.

\begin{algorithm}[H]
\caption{\textsc{AdaptiveThreshold}$(J,c)$}
\label{alg:adaptive-threshold}
\begin{algorithmic}
\Statex \textbf{Input:} jobs $J$ and parameter $c>0$.
\State Fix an arbitrary test order $j_1,\ldots,j_n$.
\State Set $N\gets0$, $N_{\mathrm{def}}\gets0$, and $Q\gets\varnothing$.
\State \textbf{For} $i=1,\ldots,n$ \textbf{do}:
\State \quad Set $x\gets n-i+1$ and $y\gets(N_{\mathrm{def}}-(1+c)N)/x$.
\State \quad Set $a\gets a_c(y)$ according to \cref{eq:opt:threshold}.
\State \quad Test $j_i$ for one unit and observe $p_{j_i}$.
\State \quad \textbf{If} $p_{j_i}\le a$, process $j_i$ immediately.
\State \quad \textbf{Otherwise}, add $j_i$ to $Q$ and set $N_{\mathrm{def}}\gets N_{\mathrm{def}}+1$.
\State \quad \textbf{If} $p_{j_i}>0$, set $N\gets N+1$.
\State Process the jobs of $Q$ in nondecreasing order of their revealed
       processing times.
\end{algorithmic}
\end{algorithm}

For $c\in[r,1/\varphi]$, let $m=m(c)\in[0,1/\varphi]$ and $u=u(c)$ be
determined by the common parametrization
\cref{eq:opt:mixed-parametrization}.
The right-hand side of the first equation is strictly increasing in $c$,
as is the left-hand side in $m$.  Consequently $m(c)$ is increasing and
$u(c)$ is decreasing.  The endpoints are
\[
 (c,m,u)=\left(\frac1\varphi,\frac1\varphi,\uDet{4}\right)
 \quad\text{and}\quad
 (c,m,u)=(r,0,\uDet{5}).
\]

\begin{theorem}[Middle and high revealing upper bounds]
\label{thm:opt:optional-middle-upper}
For every finite $u\in[\uDet{4},\infty)$, the formally specified algorithm
\textnormal{\textsc{AdaptiveThreshold}} from
\cref{alg:adaptive-threshold}
satisfies
\[
 \ALG\le
 \begin{cases}
  (1+c)\OPT+O_u(n),
     & \uDet{4}\le u\le \uDet{5},
       \quad c\text{ given by \eqref{eq:opt:mixed-parametrization}},\\[1mm]
  (1+r)\OPT+O(n),
     & u\ge \uDet{5}.
 \end{cases}
\]
The constant in the last $O(n)$ is absolute and uniform over
$u\in[\uDet{5},\infty)$.
In particular these are upper bounds on the deterministic
size-asymptotic competitive ratio.
\end{theorem}

\begin{proof}

Fix a finite cap $u$ in the stated range.

\paragraph{The algorithm.}
Use \textsc{AdaptiveThreshold}$(J,c)$ from
\cref{alg:adaptive-threshold}; on the high interval use $c=r$.
For the analysis, immediately before a test let $x$ be the number of jobs not
yet tested, $N$ the number of earlier strictly positive outcomes, and $N_{\mathrm{def}}$ the
number of those outcomes deferred to the final tail.  Thus
\begin{equation}
 y=\frac{N_{\mathrm{def}}-RN}{x}\le0
 \label{eq:opt:state-y}
\end{equation}
and the threshold is given by \cref{eq:opt:threshold}.  On the mixed interval
$c\ge r$, and hence
\[
 a_c(y)\le a_c(0)\le a_r(0)=\frac1r<u-1.
 \label{eq:opt:threshold-below-cap}
\]
The same inequality holds on the high interval, where we use $c=r$.
Thus every outcome in the capped region $[u-1,u]$ is deferred.

The state $y$ never increases.  Indeed, the numerator in
\eqref{eq:opt:state-y} changes by $0$, $-R$, or $-c$ after a zero, a
positive immediate outcome, or a positive deferred outcome, respectively,
while $x$ decreases.  This monotonicity will allow us to allocate every
pair cost when its relevant endpoint becomes known.

\paragraph{Exact pair objective and the five endpoints.}
Write
\[
 \bar p_i=\min\{p_i,u-1\}.
\]
By \cref{eq:model-effective-length,lem:offline}, the clairvoyant revealing
optimum is SPT on jobs of effective lengths
$1+\bar p_i=\min\{1+p_i,u\}$.  For two jobs tested in the order $i<j$, let
\begin{equation}
 e_{ij}=
 \begin{cases}
  (p_i-p_j)_+,
     &i\text{ is processed immediately},\\
  1+(p_j-p_i)_+,
     &i\text{ is deferred and }j\text{ is immediate},\\
  1,
     &i\text{ and }j\text{ are deferred}.
 \end{cases}
 \label{eq:opt:pair-excess}
\end{equation}
The first line applies whenever $i$ is immediate, regardless of the
action taken on $j$.  Pairwise completion-time accounting for the online
schedule gives
\begin{equation}
 \ALG
 =\sum_i(1+p_i)
   +\sum_{i<j}\bigl(1+\min\{p_i,p_j\}+e_{ij}\bigr),
 \label{eq:opt:alg-pair-decomposition}
\end{equation}
whereas the common offline pair identity \cref{eq:optional-opt-pair} gives
\begin{equation}
 \OPT
 =\sum_i(1+\bar p_i)
   +\sum_{i<j}\bigl(1+\min\{\bar p_i,\bar p_j\}\bigr).
 \label{eq:opt:opt-pair-decomposition}
\end{equation}
It is therefore natural to introduce
\begin{equation}
 \Gamma_n=\sum_{i<j}
 \left(e_{ij}+\min\{p_i,p_j\}
       -R\min\{\bar p_i,\bar p_j\}\right).
 \label{eq:opt:pair-objective}
\end{equation}
Equations \eqref{eq:opt:alg-pair-decomposition}--
\eqref{eq:opt:pair-objective} yield the exact identity
\begin{equation}
 \ALG-R\OPT
 =\Gamma_n-c\binom n2
  +\sum_i\bigl[(1+p_i)-R(1+\bar p_i)\bigr].
 \label{eq:opt:gap-decomposition}
\end{equation}
The final sum is nonpositive.  Below the cap its $i$-th summand is
$-c(1+p_i)$; on the cap it is at most $1-cu\le0$.  Thus it is enough
to control $\Gamma_n$.

Fix a symbolic immediate/deferred path.  Coordinatewise convexity in
\eqref{eq:opt:pair-objective} reduces a worst outcome vector to the five
formal endpoint types
\begin{equation}
\begin{array}{c|c|c}
 \text{type}&\text{outcome}&\text{action}\\ \hline
 Z&0&\text{immediate}\\
 E&0^+&\text{immediate}\\
 M&a_c(y)&\text{immediate}\\
 Q&a_c(y)^+&\text{deferred}\\
 U&u&\text{deferred and capped}.
\end{array}
 \label{eq:opt:five-endpoints}
\end{equation}
For an immediate coordinate the two positive endpoints are $0^+$ and
the live threshold; the genuine zero is kept separate because it does not
enter $N$.  A deferred coordinate below $u-1$ moves to the live
threshold or to $u-1$.  On $[u-1,u]$, the offline contribution is flat
and the algorithmic contribution is nondecreasing, so the latter endpoint
moves to $u$.  The symbolic categories, and therefore all later
threshold states, remain unchanged during each coordinate move.  This is
the only endpoint reduction used below.

\paragraph{A local pair allocation.}
Separate the earlier $E$-outcomes from the substantive outcomes
$M,Q,U$.  Let
\[
 P=\#\{\text{earlier }M,Q,U\},\qquad
 E=\#\{\text{earlier }E\},\qquad
 N_{\mathrm{def}}=\#\{\text{earlier }Q,U\},
\]
and let $K$ be the number of earlier $U$-outcomes.  Among the ordinary
substantive outcomes, let $I,d$ be the numbers of $M,Q$-outcomes and
let $A_I,A_D$ be the sums of their processing values.  Thus
$P=I+d+K$ and $N_{\mathrm{def}}=d+K$.

At a current ordinary endpoint, with $a=a_c(y)$, allocate the relaxed
local charges
\begin{equation}
 \ell_Z=\ell_E=N_{\mathrm{def}},
 \qquad
 \ell_M=N_{\mathrm{def}}+a(x+N_{\mathrm{def}}-RP),
 \qquad
 \ell_Q=N_{\mathrm{def}}+a(N_{\mathrm{def}}-RP).
 \label{eq:opt:ordinary-local-charges}
\end{equation}
For a capped endpoint the exact charge is instead
\begin{align}
 \ell_U
 &=N_{\mathrm{def}}+A_D+uK-R\bigl(A_I+A_D+(u-1)K\bigr)
 \notag\\
 &=d-RA_I-cA_D+K h_U,
 \qquad
 h_U=1+u-R(u-1).
 \label{eq:opt:cap-local-charge}
\end{align}
These formulae follow by checking the finitely many ordered endpoint
pairs.  In particular, an earlier immediate outcome already allocated its
online processing delay when it appeared, so a current cap contributes
only $-Rp$ against that outcome.  Against an earlier ordinary deferred
outcome it contributes $p-Rp=-cp$, and against an earlier cap it
contributes $h_U$.  The same check shows
\begin{equation}
 \Gamma_n\le\sum_i\ell_i;
 \label{eq:opt:pair-allocation-dominates}
\end{equation}
the difference is exactly the sum of the nonnegative own-processing terms
allocated by the $M$-endpoints.  Formula
\eqref{eq:opt:cap-local-charge}, rather than a substitution $p=u$ into
$\ell_Q$, is essential before caps have been put into a prefix.

Threshold monotonicity gives $A_I\ge aI$ and $A_D\ge ad$.  Comparing
\eqref{eq:opt:cap-local-charge} with the hypothetical $Q$-charge at the
same state gives the exact identity
\begin{align}
 \ell_U-\ell_Q
 &=R(aI-A_I)+c(ad-A_D)
   +K\,[1-c(u-1-a)]
 \notag\\
 &\le K\,[1-c(u-1-a)].
 \label{eq:opt:cap-versus-q}
\end{align}

\paragraph{The reusable four-endpoint bank.}
The states $Z,E,M,Q$ are the four ordinary, noncap endpoints.  The potential
below proves their inequalities simultaneously for every target parameter
$c$ on the mixed interval.  The fifth endpoint $U$ is handled only through
the extra reserve.  At $c=r$ that reserve vanishes, leaving a cap-free bank
that will also support the obligatory endpoint without a second proof.
Define
\begin{equation}
 \eta=\frac{N_{\mathrm{def}}-RP}{x},
 \qquad b=\frac Ex,
 \qquad y=\eta-Rb.
 \label{eq:opt:eta-b}
\end{equation}
For $-1\le y\le0$, put
\begin{align}
 f_c(y)
 &=\frac{c+2}{4}
   +\frac{c-2y}{4}
      \left(\ln\frac{c-2y}{c+2}-1\right),
 \label{eq:opt:f-definition}\\
 h_c(y)&=a_c(y)(1+y),
 \notag\\
 G_c(y,b)&=f_c(y)+b h_c(y)+\frac R2b^2.
 \label{eq:opt:G-region-one}
\end{align}
For $y\le-1$, set instead
\begin{equation}
 G_c(y,b)=g_c(\eta),
 \qquad
 g_c(\eta)=\frac{(1+\eta)_+^2}{2R}.
 \label{eq:opt:G-region-two}
\end{equation}
The two definitions are $C^1$ across $y=-1$.  Indeed, feasibility
implies $0\le b\le1/R$ there, and both sides have value $Rb^2/2$ and
derivatives
\[
 (G_y,G_b)=(b,Rb).
\]
They are also nonnegative: $f_c(-1)=0$,
$f_c'=a_c-1\ge0$, and $h_c,b,g_c\ge0$.

Use the base potential
\begin{equation}
 W_{\mathrm{base}}(x,P,E,N_{\mathrm{def}})=xN_{\mathrm{def}}+x^2G_c(y,b).
 \label{eq:opt:base-potential}
\end{equation}
In the raw state $(x,P,E,N_{\mathrm{def}},K)$, the endpoint directions are
\begin{equation}
\begin{array}{c|ccccc}
 &x&P&E&N_{\mathrm{def}}&K\\ \hline
 v_Z&-1&0&0&0&0\\
 v_E&-1&0&1&0&0\\
 v_M&-1&1&0&0&0\\
 v_Q&-1&1&0&1&0\\
 v_U&-1&1&0&1&1.
\end{array}
 \label{eq:opt:state-directions}
\end{equation}
Let $J_T=\ell_T+\nabla W_{\mathrm{base}}\mathbin{\cdot}v_T$, where for
the moment $T\in\{Z,E,M,Q\}$.  In Region I, direct differentiation and
the identities
\begin{align}
 -2f_c+(y-R)f_c'+a_c(1+y)&=0,
 \notag\\
 1-2f_c+(y-c)f_c'+a_cy&=0
 \label{eq:opt:f-hjb-identities}
\end{align}
give
\begin{align}
 \frac{J_Z}{x}
 &=(-2f_c+yf_c')+b(-h_c+yh_c'),
 \notag\\
 \frac{J_E}{x}
 &=b[-h_c+(y-R)h_c'+R],
 \notag\\
 \frac{J_M}{x}
 &=b[-h_c+(y-R)h_c'+a_cR],
 \notag\\
 \frac{J_Q}{x}
 &=b[-h_c+(y-c)h_c'+a_cR].
 \label{eq:opt:region-one-drifts}
\end{align}
All four quantities are nonpositive.  For the first one,
\begin{equation}
 -2f_c+yf_c'=a_c(c-y)-R\le0;
 \label{eq:opt:zero-drift-sign}
\end{equation}
the right-hand side is zero at $y=-1$ and is nonincreasing because
$a_c'(c-y)-a_c\le1-a_c\le0$.  Also $h_c,h_c'\ge0$.  For the remaining
three drifts, observe that
\begin{equation}
 h_c+(c-y)h_c'
 =a_cR+(c-y)(1+y)a_c'\ge a_cR.
 \label{eq:opt:h-sign-identity}
\end{equation}
The $Q$-bracket in \eqref{eq:opt:region-one-drifts} is therefore
nonpositive; the $M$-bracket is the $Q$-bracket minus $h_c'$, and the
$E$-bracket is the $M$-bracket minus $R(a_c-1)$.

In Region II, the same calculation gives
\begin{align}
 J_Z/x=J_E/x&=-2g_c+\eta g_c',
 \notag\\
 J_M/x&=-2g_c+(\eta-R)g_c'+1+\eta,
 \notag\\
 J_Q/x&=1-2g_c+(\eta-c)g_c'+\eta.
 \label{eq:opt:region-two-drifts}
\end{align}
For $-1\le\eta\le0$, substitution of
$g_c'=(1+\eta)/R$ makes these values $-(1+\eta)/R$,
$-(1+\eta)/R$, and $0$, respectively.  For $\eta\le-1$ their signs
are immediate from $g_c=g_c'=0$.  Hence the base bank pays every
ordinary endpoint.

\paragraph{The cap reserve.}
Define
\begin{equation}
 Z(c)=2+\frac1c+\frac12\ln\!\left(1+\frac2c\right),
 \qquad
 \delta=Z(c)-u.
 \label{eq:opt:Z-delta}
\end{equation}
On the mixed curve, \eqref{eq:opt:mixed-parametrization} says precisely
that
\begin{equation}
 \delta=\operatorname{artanh}m.
 \label{eq:opt:delta-atanh}
\end{equation}
For $0\le t\le m$, set
\begin{equation}
 \beta(t)=ct\,[\delta-\operatorname{artanh}t],
 \qquad
 \mathcal V(t)=\int_t^m\beta(s)\,ds,
 \label{eq:opt:reserve-integral}
\end{equation}
and extend $\mathcal V$ by zero for $t\ge m$.  Since
$\beta(m)=0$, this extension is $C^1$.  With
\begin{equation}
 q=\frac Kx,
 \qquad
 \mu=\frac{q}{1+q},
 \qquad
 \mathcal H(q)=(1+q)^2\mathcal V(\mu),
 \label{eq:opt:reserve-H}
\end{equation}
define the full potential
\begin{equation}
 W=xN_{\mathrm{def}}+x^2\bigl(G_c(y,b)+\mathcal H(q)\bigr).
 \label{eq:opt:full-potential}
\end{equation}
It is nonnegative.

For an ordinary endpoint, $K$ is fixed while $x$ decreases, so the
reserve contributes, after division by $x$,
\begin{equation}
 L_0(q)=-2\mathcal H(q)+q\mathcal H'(q).
 \label{eq:opt:L0}
\end{equation}
A cap also increments $K$, and its contribution is
$L_0+\mathcal H'$.  On the active range $\mu\le m$, differentiation of
\eqref{eq:opt:reserve-H} gives
\begin{equation}
 \mathcal H'=2(1+q)\mathcal V-\beta(\mu),
 \qquad
 L_0=-2(1+q)\mathcal V-q\beta(\mu)\le0.
 \label{eq:opt:reserve-derivatives}
\end{equation}
Thus the reserve can only improve every ordinary drift.

For a cap, feasibility $N_{\mathrm{def}}\le P$ and $P\ge K$ gives
\begin{equation}
 y=\frac{N_{\mathrm{def}}-R(P+E)}x\le-\frac{cK}{x}=-cq.
 \label{eq:opt:y-cap-bound}
\end{equation}
While $\mu\le m$, we have $q\le m/(1-m)\le\varphi$ and hence
$cq\le1$.  Monotonicity of $a_c$, together with
\eqref{eq:opt:cap-versus-q}, yields
\begin{align}
 \frac{\ell_U-\ell_Q}{x}
 &\le cq\left[\delta-\frac12\ln(1+2q)\right]
 =:B(q).
 \label{eq:opt:cap-residual-B}
\end{align}
The change of variables in \eqref{eq:opt:reserve-H} satisfies
\begin{equation}
 \operatorname{artanh}\mu=\frac12\ln(1+2q),
 \qquad
 (1+q)\beta(\mu)=B(q).
 \label{eq:opt:B-beta}
\end{equation}
Consequently the cap residual is cancelled exactly:
\begin{equation}
 L_0+\mathcal H'+B(q)=0
 \qquad(\mu\le m).
 \label{eq:opt:reserve-cancellation}
\end{equation}
When $\mu\ge m$, the reserve vanishes.  At
$q_*=m/(1-m)$, the bracket in
\eqref{eq:opt:cap-versus-q}, evaluated at the largest possible threshold
$a_c(-cq_*)$, is zero by
\eqref{eq:opt:delta-atanh} and \eqref{eq:opt:B-beta}.  For $q\ge q_*$,
\eqref{eq:opt:y-cap-bound} only decreases the threshold, so that bracket is
nonpositive.  The fifth Hamiltonian inequality therefore holds for every
state and for arbitrary interleavings of capped and ordinary outcomes.

\paragraph{Initial calibration.}
Initially $y=b=q=0$, and
\begin{equation}
 f_c(0)=\frac12-\frac c4\ln\!\left(1+\frac2c\right).
 \label{eq:opt:f-zero}
\end{equation}
Elementary integration in \eqref{eq:opt:reserve-integral} gives
\begin{equation}
 \mathcal V(0)=\frac c2(\delta-m).
 \label{eq:opt:V-zero}
\end{equation}
Indeed,

\[
 \int_0^m t\operatorname{artanh}t\,dt
 =\frac{(m^2-1)\operatorname{artanh}m+m}{2}.
\]
Using \eqref{eq:opt:mixed-parametrization} and
\eqref{eq:opt:delta-atanh}, we obtain the exact calibration
\begin{equation}
 f_c(0)+\mathcal H(0)
 =f_c(0)+\mathcal V(0)=\frac c2.
 \label{eq:opt:initial-calibration}
\end{equation}
Thus $W_{\mathrm{initial}}=cn^2/2$.

\paragraph{Finite jobs, and uniformity on the plateau.}
We spell out the regularity needed to pass from the fluid Hamiltonian to
unit job transitions.  It is important here to restrict attention to
feasible states.  In Region I, $y\ge-1$ implies
\begin{equation}
 -xy=R(P+E)-N_{\mathrm{def}}=cP+(P-N_{\mathrm{def}})+RE\le x.
 \label{eq:opt:feasible-compactness}
\end{equation}
All terms in the middle expression are nonnegative.  Hence
$P/x,N_{\mathrm{def}}/x,q\le1/c$ and $b\le1/R$, so every normalized derivative of
the Region-I perspective in \eqref{eq:opt:base-potential} is bounded on
the feasible cone.  In Region II the raw base bank is simply
\begin{equation}
 xN_{\mathrm{def}}+\frac{(x+N_{\mathrm{def}}-RP)_+^2}{2R},
 \label{eq:opt:region-two-raw-bank}
\end{equation}
which has bounded one-sided Hessians.  The two pieces are $C^1$ at
$y=-1$, as verified above, and the positive-part square is $C^1$ at
$\eta=-1$.

For the reserve $\mathcal R(x,K)=x^2\mathcal H(K/x)$, on its active cone
$0\le q\le q_*$ one has bounded $\mathcal H''$, and
\begin{equation}
 \mathcal R_{xx}=2\mathcal H-2q\mathcal H'+q^2\mathcal H'',
 \quad
 \mathcal R_{xK}=\mathcal H'-q\mathcal H'',
 \quad
 \mathcal R_{KK}=\mathcal H''.
 \label{eq:opt:reserve-hessian}
\end{equation}
For $q\ge q_*$, $\mathcal R=0$; moreover
$\mathcal H(q_*)=\mathcal H'(q_*)=0$.  Thus the join is $C^1$ with
bounded one-sided Hessians.  At $x=0$, set $\mathcal R=0$: for $K>0$ it is
already locally zero, while at $(x,K)=(0,0)$ two-homogeneity gives zero
gradient.  Altogether $W$ is $C^{1,1}$ on the feasible state cone.  On the
mixed interval its regularity constant may depend on the fixed value of $u$.
When $c=r$ and $m=0$, however, the reserve is identically zero and the base
bank depends only on the absolute constant $r$; its regularity constant is
therefore uniform over all $u\ge \uDet{5}$.  This includes the last transition
from $x=1$ to $x=0$.

It follows that a unit transition has the Taylor remainder
\begin{equation}
 \ell_i+W_{i+1}-W_i\le O_u(1).
 \label{eq:opt:finite-taylor}
\end{equation}
The endpoint $E=0^+$ is interpreted one-sidedly.  For a fixed symbolic
word, replace all $E$-coordinates by positive values $\eps$.  Its
state trajectory is unchanged, and the exact finite pair objective, being
a finite sum of affine and minimum terms, converges as
$\eps\downarrow0$.  Equivalently, retain an
$O(\eps n^2)$ error and take this limit for each fixed $n$.
Therefore the formal $E$-Hamiltonian transfers to genuine outcome
vectors without an asymptotic loss.

Summing \eqref{eq:opt:finite-taylor}, using $W\ge0$, and applying
\eqref{eq:opt:initial-calibration} gives
\begin{equation}
 \sum_i\ell_i\le\frac c2n^2+O_u(n).
 \label{eq:opt:ell-telescope}
\end{equation}
Together with \eqref{eq:opt:gap-decomposition} and
\eqref{eq:opt:pair-allocation-dominates}, this yields
\begin{align}
 \ALG-R\OPT
 &\le -c\binom n2+\frac c2n^2+O_u(n)
 =O_u(n).
 \label{eq:opt:mixed-upper-finite}
\end{align}

At $u=\uDet{5}$, one has $c=r$, $m=0$, and
$Z(r)=\uDet{5}$.  The cap reserve is identically zero, while
\eqref{eq:opt:initial-calibration} remains valid.  If $u$ is increased
above $\uDet{5}$, the only adverse fifth-endpoint bracket
$1-r(u-1-a_r(y))$ decreases, and the capped diagonal term in
\eqref{eq:opt:gap-decomposition} also decreases.

Here is the promised uniform finite calculation.  The base bank and its
feasible cone contain no $u$, so there is one constant $C_{\rm step}$, depending
only on $r$, for which every ordinary or capped transition with $u\ge \uDet{5}$
satisfies
\[
 \ell_i+W_{i+1}-W_i\le C_{\rm step}.
\]
Because $W_0=rn^2/2$ and $W_n=0$, telescoping gives
\begin{equation}
 \sum_i\ell_i\le\frac r2n^2+C_{\rm step}n.
 \label{eq:opt:uniform-ell}
\end{equation}
In the diagonal term of \cref{eq:opt:gap-decomposition}, a value below the
cap contributes $-r(1+p_i)$.  A capped value contributes at most
$1-r u\le0$.  Hence the diagonal sum is nonpositive uniformly on
$[\uDet{5},\infty)$.  Combining
\cref{eq:opt:gap-decomposition,eq:opt:pair-allocation-dominates,eq:opt:uniform-ell}
now yields
\begin{align}
 \ALG-(1+r)\OPT
 &\le-r\binom n2+\frac r2n^2+C_{\rm step}n
   =\left(C_{\rm step}+\frac r2\right)n.
 \label{eq:opt:uniform-high-gap}
\end{align}
Thus, with $B=C_{\rm step}+r/2$,
\begin{equation}
 \ALG\le(1+r)\OPT+B n
 \qquad\text{for every finite }u\in[\uDet{5},\infty),
 \label{eq:opt:uniform-high-upper}
\end{equation}
and $B$ is independent of $u$.  The estimate is genuinely uniform rather
than the result of sending a parameter-dependent remainder to infinity.  This proves
\cref{thm:opt:optional-middle-upper}.  Finally, because $u>1$ on the
present range, every effective offline length is at least one and
$\OPT\ge n(n+1)/2$.  The additive term in
\eqref{eq:opt:mixed-upper-finite} vanishes in the size-asymptotic ratio.

\end{proof}


\subsection{Branchwise lower bounds}
\label{sec:det-branchwise-lower}

\subsubsection{The binary stopping game and its lower envelope}
\label{sec:optional-lower}

The first four branches already occur on binary instances
$p_j\in\{0,u\}$.  The stopping point must be hidden from the algorithm;
announcing it in advance gives a strictly weaker lower bound.

\begin{lemma}[Raw-safe hidden stopping]
\label{lem:opt:hidden-stopping}
Fix $u>1$, a target $R\le u$, and a stopping fraction $\alpha\in(0,1)$.
Against a deterministic algorithm, there is a legal binary adversary whose
leading excess at its first crossing is
\begin{equation}
 (u-R)(1-\sigma^2)+\sigma^2 f_{u,R}(y,b),
 \label{eq:opt:stopping-decomposition}
\end{equation}
where $0\le b\le y\le1$ and
\begin{equation}
 0\le y-\alpha<\frac1{n-v}.
 \label{eq:opt:stopping-overshoot}
\end{equation}
Moreover,
\begin{equation}
 f_{u,R}(y,b)=
 1-R+2y+(u-1)(1-R)y^2
 +b^2+2b(u-1-uy).
 \label{eq:opt:stopping-f}
\end{equation}
Here $\sigma\in[0,1]$ is the fraction of jobs not already executed raw.
Consequently, if
\[
 \min_{0\le b\le\alpha}f_{u,R}(\alpha,b)\ge0,
\]
then the adversary gives ratio at least $R-O_u(1/n)$.
\end{lemma}

\begin{proof}
Start in a long phase.  A fresh job that is tested receives $p=u$, while a
fresh job run raw is assigned the hidden value zero.  If $v$ jobs have been
run raw and $L$ tested-long jobs have been seen, stop at the first touch for
which
\begin{equation}
 L\ge\alpha(n-v).
 \label{eq:opt:stopping-line}
\end{equation}
All untouched jobs are then assigned value zero.  Only for the purpose of
lower-bounding the remaining cost, we grant the algorithm this future
information and allow an optimal continuation; in the actual interaction,
values are still revealed only by tests.  If the line is never crossed, every job
was run raw: if even one job had been tested, then after the last first touch
the tested fraction would be one.  On the all-zero hidden instance the ratio
is $u\ge R$.

At a crossing, let $e$ of the $L$ long jobs already be complete, let $d=L-e$
be pending, and let $x=n-v-e-d$ jobs remain untouched.  The following
favorable canonical schedule lower-bounds the algorithm: run the $v$ raw
jobs; test and immediately process the $e$ completed long jobs; perform the
$d$ noncompleting long tests; test the $x$ zero jobs; and finally process the
$d$ pending long jobs.  Indeed, if a noncompleting unit test is immediately
before an available length-$u$ completion while $h$ jobs remain, the two
orders contribute $h+hu$ and $hu+(h-1)$, respectively; moving the
completion left saves one.  Raw completions are available from time zero,
and identical long jobs may be relabeled so that the $e$ completed ones are
tested first.  After the switch, a zero test-completion has length one and
therefore precedes a pending operation of length $u$.  These exchanges
respect precedence.

The exact costs of this relaxed schedule and of OPT are
\begin{align}
 A_{\rm ex}={}&
 u\left(vn-\frac{v(v-1)}2\right)
 +(1+u)\left(e(n-v)-\frac{e(e-1)}2\right) \notag\\
 &+d(n-v-e)+x(d+x)-\frac{x(x-1)}2
 +u\frac{d(d+1)}2,
 \label{eq:opt:stopping-Aexact}\\
 O_{\rm ex}={}&\frac{n(n+1)}2
 +(u-1)\frac{(e+d)(e+d+1)}2.
 \label{eq:opt:stopping-Oexact}
\end{align}
Put $\nu=v/n$, $\delta=(v+e+d)/n$, and $\lambda=e/n$.  Twice the quadratic
coefficients in \cref{eq:opt:stopping-Aexact,eq:opt:stopping-Oexact} are
\begin{align*}
 \mathcal A={}&1+2\delta(1-u\nu)+(u-1)\delta^2
 +2\nu(\nu+u-2)\\
 &+\lambda^2+2\lambda(\nu+u-1-u\delta),\\
 \mathcal O={}&1+(u-1)(\delta-\nu)^2.
\end{align*}
Now put $\sigma=1-\nu$, $y=(e+d)/(n-v)$, and $b=e/(n-v)$.
Direct expansion gives exactly \cref{eq:opt:stopping-decomposition} after
normalizing by $n^2/2$.  The first-crossing overshoot in
\cref{eq:opt:stopping-line} is less than one: a test increases
$e+d-\alpha(n-v)$ by one and a raw first touch increases it by $\alpha$.
This proves \cref{eq:opt:stopping-overshoot}, and in particular
$\sigma(y-\alpha)<1/n$.  There is one small domain issue: the actual
$b\le y$ can exceed $\alpha$.  Put $\bar b=\min\{b,\alpha\}$.  Then
$0\le b-\bar b\le y-\alpha$.  Since $f_{u,R}$ is Lipschitz on the unit
square,
\begin{equation}
 \sigma^2 f_{u,R}(y,b)
 \ge \sigma^2 f_{u,R}(\alpha,\bar b)-O_u(1/n).
\end{equation}
Thus the stated minimum over $0\le\bar b\le\alpha$ suffices uniformly even
when $n-v$ is small.  The discarded diagonal terms are $O_u(n)$ in the
unnormalized costs.
Finally, the adaptive transcript defines a fixed instance by replay: every
tested answer was fixed when revealed, and every raw job's zero value stayed
hidden.  This proves the lemma.
\end{proof}

\begin{proposition}[Binary lower curve]
\label{prop:opt:binary-lower}
Every deterministic algorithm has size-asymptotic ratio at least
\begin{equation}
 B(u)=
 \begin{cases}
  1,&0<u\le1,\\
  u,&\uDet{1}\le u\le \uDet{2},\\
  \Rquad(u),&\uDet{2}\le u\le \uDet{3},\\[1mm]
  \displaystyle1+\frac1{\sqrt{u-1}},&u\ge \uDet{3}.
 \end{cases}
 \label{eq:opt:binary-curve}
\end{equation}
The lower bound uses only $p_j\in\{0,u\}$.
\end{proposition}

\begin{proof}
\medskip
\noindent\textbf{Claim 1 (Two quadratic branches).}
Minimizing \cref{eq:opt:stopping-f} over the algorithm's completion choice
gives
\begin{equation}
 b_*(y)=\max\{0,uy-(u-1)\}.
 \label{eq:opt:bstar}
\end{equation}
The upper constraint $b\le y$ never binds, because
$uy-(u-1)\le y$ for $y\le1$.
There are therefore two elementary quadratic branches.  On the positive
branch, put
\begin{equation}
 A_u=u^2-u+1,
 \qquad
 D_{u,R}=A_u+(u-1)R.
\end{equation}
Then
\begin{equation}
 \min_b f_{u,R}(y,b)
 =1-R-(u-1)^2+2A_uy-D_{u,R}y^2.
 \label{eq:opt:stopping-positive-branch}
\end{equation}
Its maximizer is $\alpha=A_u/D_{u,R}$, and tangency to zero is
\begin{equation}
 \bigl(1-R-(u-1)^2\bigr)D_{u,R}+A_u^2=0.
 \label{eq:opt:stopping-tangency}
\end{equation}
Expanding \cref{eq:opt:stopping-tangency} gives
\cref{eq:opt:Rquad-equation}.  On the zero-$b$ branch,
\begin{equation}
 \min_b f_{u,R}(y,b)
 =1-R+2y+(u-1)(1-R)y^2,
\end{equation}
and tangency gives
\begin{equation}
 R=1+\frac1{\sqrt{u-1}},
 \qquad
 \alpha=\frac1{\sqrt{u-1}}.
 \label{eq:opt:zero-b-tangency}
\end{equation}
The latter point belongs to the zero-$b$ branch iff
$(u-1)^3\ge u^2$, with equality at $u=\uDet{3}$.

\medskip
\noindent\textbf{Claim 2 (The first transition point).}
For $R=u$, the maximum in
\cref{eq:opt:stopping-positive-branch} is nonnegative exactly while
\begin{equation}
 [u(u-1)]^2-u(u-1)-1\le0,
\end{equation}
whose positive equality is $u=\uDet{2}$.

\medskip
\noindent\textbf{Claim 3 (Attainment of the envelope).}
For $\uDet{1}<u\le \uDet{2}$, use $R=u$ and the maximizing point of
\cref{eq:opt:stopping-positive-branch}.  For
$\uDet{2}\le u\le \uDet{3}$, use the positive root $\Rquad(u)$ and
$\alpha=A_u/D_{u,\Rquad(u)}$; the tangency identity makes the maximum zero.
The minimizing completion mass is positive in the interior and becomes zero
at $u=\uDet{3}$.  For $u\ge \uDet{3}$, use the zero-$b$ tangency point from
\cref{eq:opt:zero-b-tangency}.  The branch conditions are precisely the two
inequalities checked above.  Finally, $\RdetRO(u)\ge1$ is trivial for $u\le1$.
In each nontrivial case \cref{lem:opt:hidden-stopping} supplies fixed binary
instances by replay.
\end{proof}


\subsubsection{The cap-free harmonic core}
\label{sec:lower}

This subsection develops the cap-free construction used by the last two
branches of the revealing curve.  Its notation is local to this subsection.
We first record its extremal specialization and then extract the scaled form
reused by the mixed branch.  The construction
has many equally large groups of positive jobs, revealed from longest to
shortest, followed by one large group of zero jobs.  The only design choice is
the gap between two consecutive positive processing times.  Choosing these
gaps harmonically will make all decisions of the online algorithm disappear
from the leading term.

\begin{theorem}[Extremal harmonic core]
\label{thm:lower}
For every deterministic online algorithm $\mathcal B$ and every $\eps>0$,
there are arbitrarily large fixed instances $I$ with unit tests and rational
processing times such that
\[
  \ALG_{\mathcal B}(I)\ge (\RdetOT-\eps)\OPT(I).
\]
\end{theorem}

\begin{proof}
\medskip
\noindent\textbf{Claim 1 (Fixed-input revelation experiment).}
Fix an integer $K$, a rational number $\xi>0$, and a rational slack
$\gamma\in(0,1)$; one may keep $\gamma=1/2$ throughout.  For a large scaling
integer $H$, the instance will contain $H$ jobs of each of $K$ positive types
$0,1,\ldots,K-1$ and $\xi H$ zero jobs.  We choose $H$ to clear all
denominators.

The adversary assigns processing times according to the order in which tests
are started.  The first $H$ tested jobs receive processing time $p_0$, the
next $H$ receive $p_1$, and so on; after the $K$ positive blocks, every
remaining job receives processing time zero.  Thus
\[
  p_0>p_1>\cdots>p_{K-1}>1,
\]
so the algorithm sees the positive types from longest to shortest and learns
about the zero jobs only at the end.

This adaptive description is only a convenient way to construct the input.
After running the interaction against a fixed deterministic algorithm, assign
to every job label the value revealed when that label was tested.  Replaying
the algorithm on this fixed assignment produces exactly the same transcript.
If the algorithm never tests or completes some job, its cost is infinite and
the lower bound is immediate; otherwise every job label is recorded.
Hence it is enough to lower-bound the cost in the revelation experiment.

\medskip
\noindent\textbf{Claim 2 (Shortest-first phase relaxation).}
Let $T_\ell$ be the start of the first test in positive block $\ell$, and let
$T_K$ be the start of the first test in the zero block.  Phase $\ell$ contains
the completions in $(T_\ell,T_{\ell+1}]$; a completion exactly at $T_\ell$
belongs to the preceding phase.  Each $T_\ell$ is the boundary between two
nonpreemptive operations.  Consequently, a previously tested job that is
unfinished at $T_\ell$ has not started processing and still needs its entire
processing time.

We repeatedly use the following elementary observation.  Suppose jobs with
remaining work $v_1,\ldots,v_m$ complete after some time $T$.  If their
completion order is $q_1,\ldots,q_m$, then the $g$th relative completion time
is at least $\sum_{h\le g}v_{q_h}$.  Summing these prefix inequalities shows
that their total relative completion time is minimized by arranging the
$v_i$ in nondecreasing order.  This is merely a lower bound; the online
algorithm need not actually use that order.

For normalized group masses $q_1,\ldots,q_m$ and remaining lengths
$v_1\le\cdots\le v_m$, this shortest-first bound has leading coefficient
\begin{equation}
 \mathcal S((q_g,v_g)_{g=1}^m)
 =\sum_{g=1}^m
   \left(\frac{v_gq_g^2}{2}
    +q_g\sum_{h<g}v_hq_h\right).
 \label{eq:lower-S}
\end{equation}
More precisely, the corresponding contribution for groups of size $q_gH$
is $H^2\mathcal S+O(H)$.  The quadratic term is simply the familiar sum
$v_g(1+\cdots+q_gH)$ inside one group; the cross term is the waiting caused
by all shorter groups.

We ensure that
\begin{equation}
  p_0-p_{K-1}<1.
  \label{eq:lower-span}
\end{equation}
In phase $\ell$, an unfinished job of an older type $j<\ell$ therefore has
remaining length $p_j$, whereas a type-$\ell$ job completed in its own phase
needs both its test and processing, of total length $1+p_\ell$.  Their order
in the shortest-first relaxation is
\[
  p_{\ell-1}<p_{\ell-2}<\cdots<p_0<1+p_\ell.
\]
At $T_K$, an untested zero job has remaining length $1$, so the relaxed tail
starts with all zero jobs and then uses the positive types from shortest to
longest.  These are the only two places where
\cref{eq:lower-span} is used.

\medskip
\noindent\textbf{Claim 3 (Finite accounting identity).}
We now record the online algorithm's possible behavior.  Normalize all job
counts by $H$.  Let
\begin{itemize}
\item $s_j\in[0,1]$ be the mass of type-$j$ jobs completed in their own
      phase;
\item $d_{j,\ell}\ge0$ be the mass of type-$j$ jobs completed in the later
      phase $\ell>j$; and
\item
\[
  X_{j,\ell}=s_j+\sum_{h=j+1}^{\ell-1}d_{j,h}\le1
  \qquad(\ell>j)
\]
be the mass of type $j$ completed before phase $\ell$.  Thus
$X_{j,\ell+1}$ is the mass completed by the end of phase $\ell$.
\end{itemize}
At time $T_\ell$, the machine has already performed the tests of the $\ell H$
jobs in earlier positive blocks and has processed every earlier job that has
completed.  Thus
\begin{equation}
  \frac{T_\ell}{H}
  \ge \ell+\sum_{j<\ell}p_jX_{j,\ell}.
  \label{eq:lower-start}
\end{equation}
For each phase, charge its completions their common start time from
\cref{eq:lower-start}, and charge their relative completion times using
\cref{eq:lower-S}.  Do the same at $T_K$ for the remaining positive jobs
and the $\xi H$ zero jobs.  Expanding and collecting equal variables gives
\begin{align}
 \frac{\ALG}{H^2}
 \ge{}& C_K
 +\sum_{j=0}^{K-1}\left(\frac{s_j^2}{2}-\theta_js_j\right)
 \nonumber\\
 &+\sum_{j<\ell}d_{j,\ell}
 \left[
   \ell-j-\theta_j
   -\sum_{j<k<\ell}(p_j-p_k)X_{k,\ell+1}
 \right]-O(1/H),
 \label{eq:lower-accounting}
\end{align}
where $K,\xi,\gamma$ are fixed while $H\to\infty$, and
\begin{align}
 \theta_j
   &=1-\xi(p_j-1)-\sum_{k>j}(p_j-p_k-1),
   \label{eq:lower-theta}\\
 C_K
   &=\frac{\xi^2}{2}+2\xi K+K^2+P_K,
 &
 P_K&=\sum_{j=0}^{K-1}\left(j+\frac12\right)p_j.
 \label{eq:lower-CP}
\end{align}

For completeness, here is how to audit
\cref{eq:lower-accounting} without guessing its origin.  In phase $\ell$,
the masses in the shortest-first relaxation are
\[
 d_{\ell-1,\ell},d_{\ell-2,\ell},\ldots,d_{0,\ell},s_\ell
\]
with respective lengths
\[
 p_{\ell-1},p_{\ell-2},\ldots,p_0,1+p_\ell.
\]
Multiply their total mass by the right-hand side of
\cref{eq:lower-start} and add the functional
\cref{eq:lower-S}.  In the final phase use masses
\[
 \xi,\ 1-X_{K-1,K},\ldots,1-X_{0,K}
\]
and lengths $1,p_{K-1},\ldots,p_0$.  The constant terms sum to $C_K$.
The terms containing only $s_j$ are $s_j^2/2-\theta_js_j$; the coefficient
of each $d_{j,\ell}$ is precisely the bracket in
\cref{eq:lower-accounting}.  Thus the displayed identity is just the
phasewise shortest-first bound with the phase start times retained.  No
optimality assumption about the online schedule is hidden in it.

The meaning of the expression is useful.  Completing an $s_j$ fraction of a
type in its own phase creates a quadratic completion cost but saves later
waiting cost, producing $s_j^2/2-\theta_js_j$.  A completion postponed to a
later positive phase is represented by $d_{j,\ell}$ and its bracket.  We next
choose the processing times so that every such bracket is nonnegative and all
$\theta_j$ are equal.

\medskip
\noindent\textbf{Claim 4 (Harmonic cancellation).}
Set
\begin{equation}
  p_{K-1}=1+\frac{\gamma}{\xi},
  \qquad
  p_i=p_{i+1}+\frac1{\xi+K-1-i}
  \quad(0\le i<K-1).
  \label{eq:lower-harmonic}
\end{equation}
In other words, measured upward from the shortest type, the consecutive gaps
are
\[
  \frac1{\xi+1},\frac1{\xi+2},\ldots,\frac1{\xi+K-1}.
\]
This is the key design choice.  If $m=K-1-j$, then
\begin{align*}
 p_j-1
   &=\frac\gamma \xi+\sum_{h=1}^m\frac1{\xi+h},\\
 \sum_{k>j}(p_j-p_k)
   &=\sum_{h=1}^m\frac{h}{\xi+h}.
\end{align*}
Adding the first line multiplied by $\xi$ to the second line makes each
denominator cancel:
\begin{equation}
 \xi(p_j-1)+\sum_{k>j}(p_j-p_k)
 =\gamma+\sum_{h=1}^m\frac{\xi+h}{\xi+h}
 =\gamma+m.
 \label{eq:lower-telescope}
\end{equation}
Substitution in \cref{eq:lower-theta} gives the same value for every type,
\begin{equation}
  \theta_j=1-\gamma.
  \label{eq:lower-theta-constant}
\end{equation}
This is why the increments in \cref{eq:lower-harmonic} are harmonic: the
factor $\xi$ coming from the zero block and the multiplicity $h$ coming from the
$h$ shorter positive types add up to the denominator $\xi+h$.

The delayed-completion terms are now harmless as well.  Since
$X_{k,\ell+1}\le1$, their brackets are at least
\begin{align*}
 &\ell-j-(1-\gamma)
   -\sum_{j<k<\ell}(p_j-p_k)\\
 &\hspace{20mm}
 =\gamma+\sum_{j<k<\ell}\bigl(1-(p_j-p_k)\bigr)>0,
\end{align*}
where the final inequality follows from
\cref{eq:lower-span}.  We may therefore drop every term containing
$d_{j,\ell}$.  The remaining minimization is not a multivariable optimization
problem: complete one square for each $j$,
\[
  \frac{s_j^2}{2}-(1-\gamma)s_j
  =\frac12\bigl(s_j-(1-\gamma)\bigr)^2
   -\frac12(1-\gamma)^2.
\]
Because $1-\gamma\in(0,1)$, this minimum is feasible for
$s_j\in[0,1]$.  Hence every online algorithm obeys
\begin{equation}
 \liminf_{H\to\infty}\frac{\ALG}{H^2}
 \ge Q_K
 \defeq C_K-\frac K2(1-\gamma)^2.
 \label{eq:lower-QK}
\end{equation}

It remains to compute the offline benchmark.  Knowing all processing times,
the optimum runs the zero jobs first and the positive types from shortest to
longest.  Equivalently, the pairwise contribution of two jobs is
$1+\min\{p_i,p_j\}$.  Its leading coefficient is
\begin{equation}
  D_K=\frac{\xi^2}{2}+\xi K+\frac{K^2}{2}+P_K,
  \qquad
  \lim_{H\to\infty}\frac{\OPT}{H^2}=D_K.
  \label{eq:lower-DK}
\end{equation}
The four terms respectively account for zero--zero pairs, zero--positive
pairs, the unit-test part of positive--positive pairs, and the processing
part of positive--positive pairs.  Thus the finite-$K$ lower-bound
certificate is simply $Q_K/D_K$.

\medskip
\noindent\textbf{Claim 5 (Smooth limit and its maximum).}
Write
\[
  \alpha=\frac K\xi
\]
for the ratio of positive mass to zero mass, and keep $\alpha$ fixed while
$K\to\infty$.  We restrict to $0<\alpha<e-1$.  Then
\[
 p_0-p_{K-1}
 =\sum_{h=1}^{K-1}\frac1{\xi+h}
 <\ln\!\left(1+\frac K\xi\right)
 =\ln(1+\alpha)<1,
\]
which verifies \cref{eq:lower-span}.

The processing contribution $P_K$ has the exact form
\begin{equation}
 P_K
 =\frac{K^2}{2}\left(1+\frac\gamma \xi\right)
  +\frac12\sum_{h=1}^{K-1}\frac{(K-h)^2}{\xi+h}.
 \label{eq:lower-P-exact}
\end{equation}
Indeed, the common baseline $1+\gamma/\xi$ contributes the first term, and
swapping the order of summation shows that the increment $1/(\xi+h)$ receives
total weight $(K-h)^2/2$.  After division by $\xi^2$, the sum becomes a
Riemann integral:
\begin{align}
 \frac{P_K}{\xi^2}&\longrightarrow I(\alpha)
   =\frac{\alpha^2}{2}
    +\frac12\int_0^\alpha\frac{(\alpha-t)^2}{1+t}\,dt
   \nonumber\\
 &=\int_0^\alpha(\alpha-t)\bigl(1+\ln(1+t)\bigr)\,dt
   \nonumber\\
 &=\frac12(1+\alpha)^2\ln(1+\alpha)
   -\frac\alpha2-\frac{\alpha^2}{4}.
 \label{eq:lower-I}
\end{align}
The second integral form follows either by integration by parts or by summing
the limiting processing-time curve $1+\ln(1+t)$ against the remaining
positive mass $\alpha-t$.
The subtraction in \cref{eq:lower-QK} is only $O(K)$, whereas
$C_K,D_K=\Theta(K^2)$.  Since $C_K-D_K=\xi K+K^2/2$, we obtain
\begin{equation}
 R(\alpha)
 \defeq\lim_{K\to\infty}\frac{Q_K}{D_K}
 =1+
 \frac{\alpha+\alpha^2/2}
 {\frac12+\frac\alpha2+\frac{\alpha^2}{4}
  +\frac12(1+\alpha)^2\ln(1+\alpha)}.
 \label{eq:lower-Ralpha}
\end{equation}

This last one-variable maximization is where the constant in the theorem
comes from.  Put $z=(1+\alpha)^2$.  Then
\begin{equation}
  R(z)=1+\frac{2(z-1)}{1+z+z\ln z}.
  \label{eq:lower-Rz}
\end{equation}
The derivative of the fraction has the sign of
\[
  3-z+\ln z.
\]
For $z>1$ this expression is strictly decreasing, is positive at $z=1$, and
tends to $-\infty$.  Hence it has one zero $\uDet{5}$, characterized by
\[
  \uDet{5}-3=\ln \uDet{5}.
\]
The derivative is positive before $\uDet{5}$ and negative afterward, so this is
the global maximum.  At the root,
\[
  1+\uDet{5}+\uDet{5}\ln \uDet{5}=(\uDet{5}-1)^2,
\]
and therefore
\[
  R(\uDet{5})=1+\frac{2}{\uDet{5}-1}=\RdetOT.
\]
Moreover, $\sqrt{\uDet{5}}-1=1.12255\ldots<e-1$, so the maximizer lies in the
range where \cref{eq:lower-span} holds.

\medskip
\noindent\textit{Final parameter choice.}
Choose a rational $\alpha$ sufficiently close to
$\alpha_\star=\sqrt{\uDet{5}}-1$, and choose any rational
$\gamma\in(0,1)$.  For a sufficiently large $K$, set $\xi=K/\alpha$ and use
\cref{eq:lower-harmonic}; then $Q_K/D_K>\RdetOT-\eps/2$.
All parameters are rational.  Let $H$ tend to infinity through multiples of
their common denominators.  \Cref{eq:lower-QK,eq:lower-DK} make the actual ratio exceed
$\RdetOT-\eps$ for all sufficiently large $H$.  Finally, the replay argument
at the beginning of the proof converts the revelation experiment into a fixed
instance for the given deterministic algorithm.
\end{proof}

The mixed branch needs the same construction at an arbitrary total mass and
at a prescribed point of its limiting curve.  The next lemma packages this
rescaled form so that no phase accounting has to be repeated later.

\begin{lemma}[Scaled harmonic core]
\label{lem:opt:harmonic-block}
Fix a total normalized mass $s>0$ and a parameter $t\in[1,e)$.  Against
every deterministic obligatory-testing algorithm, there are scales
$M_\nu\to\infty$ and finite rational multilevel instances with
$sM_\nu+o(M_\nu)$ jobs such that
\begin{equation}
 \frac{\OPT}{M_\nu^2}\longrightarrow
 O_0=\frac{s^2}{4}(1+t^{-2}+2\ln t),
 \qquad
 \liminf_{\nu\to\infty}\frac{\ALG}{M_\nu^2}
 \ge A_0^{\rm lb}
 \defeq O_0+\frac{s^2}{2}(1-t^{-2}).
 \label{eq:opt:harmonic-coefficients}
\end{equation}
Their normalized total positive processing mass converges to
\begin{equation}
 \frac1{M_\nu}\sum_{j:p_j>0}p_j\longrightarrow L=s\ln t,
 \label{eq:opt:harmonic-processing-mass}
\end{equation}
and their largest processing time converges to $1+\ln t<2$.
The adaptive revelation is converted by replay into a fixed instance.
\end{lemma}

\begin{proof}
The case $t=1$ is obtained by continuity from instances with vanishing
positive mass, so suppose $1<t<e$.  In the construction above set
$\alpha=t-1=K/\xi$.  Before rescaling, the total normalized mass is
$K+\xi=\xi t$.  Give every unit of normalized mass the common factor
$s/(\xi t)$.

The phase accounting, harmonic cancellation, and completion of squares in
\cref{eq:lower-accounting,eq:lower-telescope,eq:lower-QK} are homogeneous of
degree two in the job masses.  The Riemann-sum limit
\cref{eq:lower-I} substituted into \cref{eq:lower-DK} therefore gives
\[
 O_0=\frac{s^2}{4}(1+t^{-2}+2\ln t).
\]
Likewise $C_K-D_K=\xi K+K^2/2$, while the $O(K)$ square-completion loss in
\cref{eq:lower-QK} disappears after division by the quadratic scale.  After
the same rescaling this gives the guaranteed lower benchmark
\[
 A_0^{\rm lb}-O_0=\frac{s^2}{2}(1-t^{-2}).
\]
Finally, summing the limiting processing-time curve
$1+\ln(1+x)$ over the positive mass gives $L=s\ln t$, and
\cref{eq:lower-harmonic} gives $p_0\to1+\ln t$.  Rational approximation,
clearing denominators, and the replay argument used in the theorem produce a
diagonal sequence of arbitrarily large fixed rational instances with the
stated limit and liminf bounds.
\end{proof}

\subsubsection{The unrestricted mixed lower bound}
\label{sec:opt:mixed-lower}

For $\uDet{4}<u<\uDet{5}$, neither the binary family nor the pure cap-free
family is tight.  The correct adversary puts a capped block before a scaled
harmonic core.  Two lemmas ensure that raw
execution and early processing of caps cannot escape it.

The cap-free construction needed below is exactly the scaled harmonic core
from \cref{lem:opt:harmonic-block}; no second phase-accounting proof is
needed here.  We now prepend capped jobs to that reusable core.

\begin{lemma}[Cap deferral]
\label{lem:opt:cap-deferral}
Prepend capped mass $m$ with value $p=u$ to a later cap-free block of mass
$s=1-m$ and total processing mass $L$.  Suppose every later value satisfies
$1+p\le u$ and
\begin{equation}
 s(u-1)-L-m\ge0.
 \label{eq:opt:cap-deferral-condition}
\end{equation}
Then processing any capped job before all later jobs are complete cannot
decrease the leading online cost.
\end{lemma}

\begin{proof}
A cap completed after its test block but before a later job of value $p$
has mutual cost $1+u$; leaving it pending until after that job gives $2+p$.
The exchange saves $u-1-p\ge0$.  If cap mass $q$ is instead completed inside
the cap-test block, its pairs with all later jobs add
$q[s(u-1)-L]$.  Its largest possible saving against not-yet-tested caps is
$qm-q^2/2$.  Relative to leaving every cap pending, the change is therefore
at least
\begin{equation}
 q[s(u-1)-L-m]+\frac{q^2}{2}\ge0.
\end{equation}
The finite diagonal discrepancy is lower order.
\end{proof}

\begin{lemma}[Capped-block accounting]
\label{lem:opt:mixed-accounting}
After cap deferral, the offline coefficient $O$ and the guaranteed online
lower benchmark $A^{\rm lb}$ satisfy
\begin{align}
 A^{\rm lb}={}&A_0^{\rm lb}+m(1-m)+m(1+L)+\frac u2m^2,
 \label{eq:opt:mixed-A}\\
 O={}&O_0+m(1-m+L)+\frac u2m^2.
 \label{eq:opt:mixed-O}
\end{align}
In particular, $A^{\rm lb}-O=(A_0^{\rm lb}-O_0)+m$.
\end{lemma}

\begin{proof}
A cap--later pair contributes $2+p$ online and $1+p$ offline,
whereas a cap--cap pair contributes $2+u$ online and $u$ offline; collecting
these pairs yields the two cross terms and the cap squares above.
Subtracting the two identities gives the final assertion.
\end{proof}

\begin{lemma}[Raw-safe quota and replay]
\label{lem:opt:mixed-raw-safe}
Fix $m\in(0,1)$ and a later cap-free block satisfying the hypotheses of
\cref{lem:opt:cap-deferral}.  Against every deterministic revealing algorithm
there are fixed instances whose asymptotic excess is at least the cap-free
hybrid benchmark $A^{\rm lb}-R\,O$ from
\cref{eq:opt:mixed-A,eq:opt:mixed-O} for every target
$R\le u$.
\end{lemma}

\begin{proof}
In the initial phase, a tested fresh job receives $p=u$, while a job run raw
is assigned hidden value zero.  If $v$ raw jobs and $\ell$ tested caps have
been touched, stop at the first crossing
\begin{equation}
 \frac{\ell}{n-v}\ge m.
 \label{eq:opt:mixed-quota}
\end{equation}
With $S=n-v$, the overshoot is less than one, so $\ell=mS+O(1)$.  If no
crossing occurs, no test can have occurred: after all first touches, any
positive $\ell$ would make the ratio in \cref{eq:opt:mixed-quota} equal one.
Thus every job was run raw and the all-zero hidden instance has ratio
$u\ge R$.

At a crossing, scale the prescribed later instance to the remaining
$S-\ell$ first touches, absorbing the one-job cap overshoot and integer
rounding into its zero group.  A later test reveals the next value in this
fixed type stream; a later raw action consumes the next value but keeps it
hidden from the revealing algorithm.

Move every initial raw-zero completion to the start.  It has the same length
$u$ as cap processing, and crossing a unit test only lowers completion cost.
With
\begin{equation}
 Z_v=vn-\frac{v(v-1)}2,
\end{equation}
this prefix contributes $uZ_v$ online and $Z_v$ offline.  Delete it.  The
remaining $S$ labels consist, up to one-job rounding, of capped fraction $m$
and the prescribed later cap-free block.

If the revealing algorithm later runs a fresh assigned value $p$ raw, a
physical cap-free simulation tests and immediately processes that job in
time $1+p\le u$, suppresses the answer from the virtual revealing algorithm,
and idles until its virtual $u$-block ends.  Inductively every physical
completion is no later than its virtual revealing completion.  Since the
finite-cap and cap-free offline lengths are both $1+p$ on the later block,
\begin{equation}
 \ALG_{\mathsf{RO}}-R\OPT_{\mathsf{RO}}
 \ge (u-R)Z_v+
 \bigl(\ALG_{\rm hyb}-R\OPT_{\rm hyb}\bigr)-O_u(n).
 \label{eq:opt:mixed-raw-decomposition}
\end{equation}
The first term is nonnegative.  Finally, record every assigned value on its
job label.  Replaying this fixed vector produces the same transcript because
the values learned only by the physical simulation were hidden from the
virtual algorithm.  This proves legality.
\end{proof}

\begin{proposition}[Mixed lower branch]
\label{prop:opt:mixed-lower}
For every $\uDet{4}\le u\le\uDet{5}$, every deterministic revealing
algorithm has size-asymptotic competitive ratio at least $\Rmix(u)$.
Equivalently,
\begin{equation}
 \RdetRO(u)\ge\Rmix(u).
 \label{eq:opt:mixed-lower-final}
\end{equation}
\end{proposition}

\begin{proof}
\medskip
\noindent\textbf{Claim 1 (Two-parameter lower ratio).}
Apply \cref{lem:opt:harmonic-block} with $s=1-m$.  Substitution into
\cref{eq:opt:mixed-A,eq:opt:mixed-O} gives the guaranteed two-parameter
lower ratio
\begin{equation}
 \mathcal R_u(m,t)=1+
 \frac{m+\frac{(1-m)^2}{2}(1-t^{-2})}
 {\frac{(1-m)^2}{4}(1+t^{-2}+2\ln t)
  +m(1-m)(1+\ln t)+\frac u2m^2}.
 \label{eq:opt:mixed-ratio}
\end{equation}
\medskip
\noindent\textbf{Claim 2 (Mixed parametrization).}
At an interior stationary point with value $1+c$, differentiation and the
zero-gap equation give
\begin{equation}
 \ln t=u-2-\frac1c,
 \qquad
 t^2=\frac{1-m}{1+m}\left(1+\frac2c\right).
 \label{eq:opt:mixed-stationarity}
\end{equation}
Eliminating $t$ gives precisely
\cref{eq:opt:mixed-parametrization}.  Conversely, those equations make
$\mathcal R_u(m,t)=1+c$.

\medskip
\noindent\textbf{Claim 3 (Feasibility and uniqueness).}
The branch runs from
\begin{equation}
 (u,m,t)=\left(\uDet{4},\frac1\varphi,1\right)
 \quad\text{to}\quad
 (u,m,t)=(\uDet{5},0,\sqrt{\uDet{5}}).
 \label{eq:opt:mixed-endpoints}
\end{equation}
It satisfies $u\ge\uDet{4}$, $\ln t<1$, and $m\le1/\varphi$.  Therefore
\begin{equation}
 (1-m)(u-1-\ln t)
 >(1-m)\varphi\ge\frac1\varphi\ge m,
\end{equation}
so \cref{eq:opt:cap-deferral-condition} holds; also $1+p<3<u$ for every
later positive value in a sufficiently fine finite construction.  The
functions
\begin{equation}
 g(c)=1-\frac1c+\frac12\ln\!\left(1+\frac2c\right),
 \qquad h(m)=\operatorname{artanh}m-m
\end{equation}
have $g'(c)=2/[c^2(c+2)]>0$ and
$h'(m)=m^2/(1-m^2)>0$.  Hence $m(c)=h^{-1}(g(c))$ is strictly increasing,
while
\begin{equation}
 u(c)=1+\frac2c-m(c)
\end{equation}
is strictly decreasing.  The identities $g(r)=h(0)=0$ and
$g(1/\varphi)=h(1/\varphi)$ give the two endpoints in
\cref{eq:opt:mixed-endpoints}; consequently every
$u\in[\uDet{4},\uDet{5}]$ has exactly one parameter pair.

\medskip
\noindent\textbf{Claim 4 (Fixed finite instances).}
To pass to fixed finite instances, first approximate $(m,t)$ rationally
while retaining the strict cap inequalities, then choose a sufficiently
fine harmonic family, and finally multiply all rational block sizes by a
common integer $H$.  For fixed harmonic depth, quota overshoot, diagonals,
and rounding are $O_u(H)$ whereas the displayed costs are $\Theta(H^2)$.
Send first $H$, then the harmonic depth, and finally the rational
approximation to their limits.  Together with
\cref{lem:opt:mixed-raw-safe}, this proves the proposition.
The construction above applies directly in the open interval.  At
$u=\uDet{4}$ and $u=\uDet{5}$ the same statement follows by approaching from
the interior, so the endpoints are included.
\end{proof}

\subsubsection{The plateau lower bound by cap-free transfer}

The harmonic construction contains no capped outcomes.  The next lemma
shows that any revealing algorithm at cap $u$ induces an obligatory-testing
algorithm on inputs with $p_j\le u-1$, so the cap-free lower bound transfers
without changing either the transcript comparison or the offline optimum.

\begin{lemma}[Bounded cap-free transfer]
\label{lem:opt:obligatory-transfer}
Let $u\ge1$, and let $L_{\mathsf{OT}}(P)$ denote the deterministic
size-asymptotic value in the cap-free limit, restricted to
$0\le p_j\le P$.  Then
\begin{equation}
 \RdetRO(u)\ge L_{\mathsf{OT}}(u-1).
 \label{eq:opt:obligatory-transfer}
\end{equation}
In particular,
\begin{equation}
 \RdetRO(u)\ge \RdetOT
 \qquad
 \left(u\ge1+\frac1r=2.752620747896442\ldots\right).
 \label{eq:opt:high-lower}
\end{equation}
\end{lemma}

\begin{proof}
Given a finite-cap algorithm, construct a cap-free algorithm on instances
with $p_j\le u-1$ by running a virtual copy of the finite-cap algorithm.  Copy
ordinary tests and processing actions.  When the virtual algorithm runs an
untouched job raw for $u$ time, physically test and immediately process it in
$1+p_j\le u$ time, hide the answer from the virtual copy, and wait for the
virtual clock.  Every physical completion is no later than its virtual
completion.  Moreover, $\min\{u,1+p_j\}=1+p_j$, so finite-cap and cap-free
optima coincide.  Thus we obtain the following game-value comparison.  For
every finite-cap algorithm
$\mathcal A$, the simulation produces a cap-free algorithm $\mathcal B$ and,
pointwise on every instance with $p_j\le u-1$,
\[
 \ALG_{\mathcal A}\ge\ALG_{\mathcal B},
 \qquad \OPT_u=\OPT_\infty.
\]
Consequently
\[
 \CR_\infty(\mathcal A)
 \ge \CR_\infty(\mathcal B; p\le u-1)
 \ge L_{\mathsf{OT}}(u-1).
\]
Taking the infimum over $\mathcal A$ proves
\cref{eq:opt:obligatory-transfer}.

The harmonic cap-free family attaining $\RdetOT$ has limiting largest
processing time
\begin{equation}
 1+\ln(1+y_*)
 =1+\frac12\ln \uDet{5}
 =\frac{\uDet{5}-1}{2}=\frac1r,
 \qquad y_*=\sqrt{\uDet{5}}-1.
\end{equation}
Approximation from below also gives the endpoint, proving
\cref{eq:opt:high-lower}.
\end{proof}


\subsection{Assembly of the curve}
\label{sec:det-curve-assembly}

\subsubsection{Shape of the candidate curve}

Before matching the bounds, we verify that the six displayed branches have
the advertised shape.  The only nontransparent part is the algebraic branch
$\Rquad$; the remaining monotonicity follows directly from the reciprocal
formula and the mixed parametrization.

\begin{lemma}[Monotonicity of the deterministic curve]
\label{lem:opt:curve-monotonicity}
The function $\Rquad(u)$ is strictly decreasing on
$[\uDet{2},\uDet{3}]$.  The reciprocal and mixed branches are strictly
decreasing on their nonconstant ranges.  Consequently the six-branch
function in \cref{eq:opt:full-curve} attains its global maximum only at
$u=\uDet{2}$.
\end{lemma}

\begin{proof}
Implicit differentiation of $T_s(\Rquad(u))=0$, with $s=u-1$, gives
\begin{equation}
 \Rquad'(u)=
 -\frac{R^2+(3s^2+2s)R-3(s+1)^2}
 {2sR+s^3+s^2+1}.
 \label{eq:opt:Rquad-derivative}
\end{equation}
Using $T_s(R)=0$, its numerator factors as
\[
 \frac{2s^3+s^2-1}{s}
 \left(R-
 \frac{2s^3+3s^2-2}{2s^3+s^2-1}\right).
\]
On the relevant interval $s\ge \uDet{2}-1>2/3$, the prefactor is
positive, while the fraction in parentheses is at most $5/3$ because
$4s^3-4s^2+1>0$.  The sign certificate
\eqref{eq:opt:T-five-thirds} gives $R>5/3$, and hence
$\Rquad'(u)<0$.  The reciprocal branch is visibly decreasing.  In the
mixed parametrization, $m(c)$ is increasing and
$u(c)=1+2/c-m(c)$ is decreasing, so $1+c$ decreases with $u$.
The first nonconstant branch increases up to $u=\uDet{2}$, all subsequent
nonconstant branches decrease, and the last branch is constant.  This proves
the global-maximum assertion.
\end{proof}

\subsubsection{Exact joins and completion of the proof}

We next check that adjacent formulas agree at every transition point.  Once
this consistency is recorded, the theorem follows by pairing the upper and
lower result already proved for each interval.

\begin{lemma}[Exact joins]
\label{lem:opt:curve-joins}
The six branches in \cref{eq:opt:full-curve} agree at all five transition
points.  More explicitly, the first two branches both equal one at
$\uDet{1}$, and
\begin{align}
 \Rquad(\uDet{2})&=\uDet{2},
 \notag\\
 \Rquad(\uDet{3})&=1+\frac1{\sqrt{s_0}}
              =1.682327803828019\ldots,
 \notag\\
 1+\frac1{\sqrt{\uDet{4}-1}}&=\varphi
              =\Rmix(\uDet{4}),
 \notag\\
 \Rmix(\uDet{5})&=1+r=\RdetOT.
 \label{eq:opt:curve-joins}
\end{align}
\end{lemma}

\begin{proof}
The equality at $\uDet{1}=1$ is immediate.  Substitution of the defining
equations for $\uDet{2}$ and $s_0=\uDet{3}-1$ into $T_s$ gives the first
two displayed identities; uniqueness of the positive root identifies the
values with $\Rquad$.  Since
$\uDet{4}-1=\varphi+1=\varphi^2$, the reciprocal branch equals
$1+1/\varphi=\varphi$.  The two remaining equalities are exactly the endpoint
conditions of the mixed parametrization in
\cref{eq:opt:mixed-endpoints}.
\end{proof}

\begin{proof}[Proof of \cref{thm:optional-curve}]
For the upper bounds, \cref{lem:opt:raw-upper} gives the first two branches,
\cref{lem:opt:ute-endpoint-game} gives the algebraic branch, and
\cref{lem:opt:zero-prefix} gives the reciprocal branch.  The mixed branch
and the uniform plateau follow from
\cref{thm:opt:optional-middle-upper}.

For the matching lower bounds, \cref{prop:opt:binary-lower} supplies the
first four branches, \cref{prop:opt:mixed-lower} supplies the fifth, and
\cref{lem:opt:obligatory-transfer,thm:lower} supplies the plateau with
cap-free instances.  The exact joins are given by
\cref{lem:opt:curve-joins}, and the claimed shape follows from
\cref{lem:opt:curve-monotonicity}.  The finite upper bounds have the stated
$O_u(n)$ remainders, while
\cref{thm:opt:optional-middle-upper} gives one absolute $O(n)$ remainder on
the plateau.  This proves every assertion of the theorem.
\end{proof}

\begin{remark}[Representative mixed values]
For orientation, some values on the mixed branch are
\begin{center}
\begin{tabular}{@{}rrrr@{}}
\toprule
$u$ & $\RdetRO(u)$ & $m$ & $t$\\
\midrule
$3.800000$ & $1.598851588$ & $0.539725635$ & $1.138984621$\\
$4.000000$ & $1.583283633$ & $0.428863572$ & $1.330517832$\\
$4.071679486$ & $1.579204806$ & $0.381330539$ & $1.412236296$\\
$4.300000$ & $1.571692026$ & $0.198387083$ & $1.734651379$\\
\bottomrule
\end{tabular}
\end{center}
\end{remark}

\begin{remark}[Asymptotic scope]
\label{rem:opt:curve-audit}
All statements above concern the fixed-$u$ size-asymptotic ratio.  Below the
plateau, the upper bounds have additive error $O_u(n)$.  On
$u\in[\uDet{5},\infty)$ the same algorithm has one uniform $O(n)$ error.
The unbounded obligatory endpoint is not included in the statement of the
curve; it is derived in \cref{sec:obligatory}.
\end{remark}


\subsection{The obligatory endpoint}
\label{sec:obligatory}

The obligatory model is the unbounded endpoint of revealing optimization,
but the passage to that endpoint is not a formal substitution
$u=\infty$: both the algorithm and the additive error must be uniform in
$u$.  The finite curve has just established exactly those two properties,
so the endpoint now follows by a short transfer.

\begin{theorem}[Deterministic obligatory testing]
\label{thm:det-obligatory}
Let $\uDet{5}>1$ be the unique solution of
$\uDet{5}-3=\ln \uDet{5}$, and set
\[
 \RdetOT=1+\frac{2}{\uDet{5}-1}
 =1.570574096649395\ldots .
\]
For every deterministic online algorithm $\mathcal B$ and every $\eps>0$,
there are arbitrarily large obligatory-testing instances $I$ such that
\[
 \ALG_{\mathcal B}(I)\ge(\RdetOT-\eps)\OPT(I).
\]
Conversely, there are a deterministic algorithm $\mathcal A$ and an absolute
constant $B$ such that every $n$-job obligatory instance satisfies
\[
 \ALG_{\mathcal A}(I)\le\RdetOT\OPT(I)+Bn.
\]
Consequently the optimal deterministic size-asymptotic competitive ratio is
exactly $\RdetOT$.
\end{theorem}

\begin{proof}
For the upper bound, use \textsc{AdaptiveThreshold} with $c=r$.  This
algorithm never runs a job raw and is independent of $u$.  Given a finite
obligatory instance, choose any finite
\[
 u>\max\{\uDet{5},1+\max_j p_j\}.
\]
Its execution is identical in the revealing model with cap $u$, and every
offline effective length is
$\min\{u,1+p_j\}=1+p_j$.  Thus the revealing and obligatory optima coincide.
The constant $B$ in \cref{eq:opt:uniform-high-upper} is independent of $u$,
so the same inequality holds for the obligatory instance.

For the lower bound, regard an arbitrary obligatory algorithm $\mathcal B$
as a revealing algorithm at the fixed cap $u=\uDet{5}$ that simply never uses raw
execution.  The plateau lower bound in \cref{thm:optional-curve} may be
witnessed by cap-free instances with $p_j<\uDet{5}-1$.  On those instances the
online transcript is unchanged and
$\min\{\uDet{5},1+p_j\}=1+p_j$, so again the two offline optima coincide.  The
finite-cap lower inequality therefore gives the claimed obligatory lower
bound for arbitrarily large instances.
\end{proof}

